% GENERATED FILE. Edit canonical lessons, then run npm run generate:books.
% canonical-source-hash: fnv1a64:91e1b2b25c26e4ce
% canonical-lessons: JA-C32-sakana, JA-C32-niku, JA-C32-tamago, JA-C32-yasai, JA-C32-kudamono

\chapter{Fish, Meat, Egg, Vegetables, Fruit}
\label{ch:ja-fish-meat-egg-vegetables-fruit}
\hlchaptermodality{32}

\begin{chapteropening}
\textbf{By the end of this chapter:} \emph{I can say a fish, meat, an egg, vegetables and fruit.}

Complete the last lesson of chapter 32: I can say a fish, meat, an egg, vegetables and fruit.
\end{chapteropening}

\section[\texorpdfstring{\ja{さかな} (sakana)}{さかな (sakana)}]{\ja{さかな} (sakana) — a fish}

\label{lesson:JA-C32-sakana}



\begin{quote}
Before the new one: say the Japanese for to move, then the Japanese for rice; a meal.
\end{quote}

\subsection*{You'll want to know: \ja{さかな}}
\textbf{\ja{さかな}} — \emph{sakana} — \textquotedblleft{}a fish\textquotedblright{}.

\begin{cousinweb}
\textbf{\ja{さかな}} (\emph{sakana}) is a native Japanese word, inherited rather than borrowed from Chinese, and its older origin is not settled. Its kanji was attached to it later as a kun'yomi, a native reading given to a Chinese character (JMdict, KANJIDIC2).
\end{cousinweb}

\begin{grammarlens}[title={What you've built}]
One more word for this chapter.
\end{grammarlens}

\subsection*{Your turn}
\begin{itemize}
\raggedright
  \item \emph{Say it:} \emph{sakana}
  \item \emph{Say it:} \emph{sakana}, once more
  \item \emph{Say it:} say \emph{sakana}
  \item \emph{Recall:} say the Japanese for to move, then the Japanese for rice; a meal, then say \emph{sakana} again
\end{itemize}

\begin{tcolorbox}[breakable,colback=teal!4,colframe=teal!35!black,title={Before you move on}]
What does \ja{さかな} mean? (\textquotedblleft{}A fish\textquotedblright{}.) Say it once more.
\end{tcolorbox}
\section[\texorpdfstring{\ja{にく} (niku)}{にく (niku)}]{\ja{にく} (niku) — meat}

\label{lesson:JA-C32-niku}



\begin{quote}
Before the new one: say the Japanese for rice; a meal, then the Japanese for a fish.
\end{quote}

\subsection*{You'll want to know: \ja{にく}}
\textbf{\ja{にく}} — \emph{niku} — \textquotedblleft{}meat\textquotedblright{}.

\begin{cousinweb}
\textbf{\ja{にく}} (\emph{niku}) is Sino-Japanese (kango): the on'yomi of a character meaning \textquotedblleft{}flesh, meat\textquotedblright{}, a reading borrowed from Chinese (Middle Chinese \emph{nyuwk}; Unihan, KANJIDIC2, Baxter–Sagart).
\end{cousinweb}

\begin{grammarlens}[title={What you've built}]
One more word for this chapter.
\end{grammarlens}

\subsection*{Your turn}
\begin{itemize}
\raggedright
  \item \emph{Say it:} \emph{niku}
  \item \emph{Say it:} \emph{niku}, once more
  \item \emph{Say it:} say \emph{niku}
  \item \emph{Recall:} say the Japanese for rice; a meal, then the Japanese for a fish, then say \emph{niku} again
\end{itemize}

\begin{tcolorbox}[breakable,colback=teal!4,colframe=teal!35!black,title={Before you move on}]
What does \ja{にく} mean? (\textquotedblleft{}Meat\textquotedblright{}.) Say it once more.
\end{tcolorbox}
\section[\texorpdfstring{\ja{たまご} (tamago)}{たまご (tamago)}]{\ja{たまご} (tamago) — an egg}

\label{lesson:JA-C32-tamago}



\begin{quote}
Before the new one: say the Japanese for a fish, then the Japanese for meat.
\end{quote}

\subsection*{You'll want to know: \ja{たまご}}
\textbf{\ja{たまご}} — \emph{tamago} — \textquotedblleft{}an egg\textquotedblright{}.

\begin{cousinweb}
\textbf{\ja{たまご}} (\emph{tamago}) is a native word, inherited rather than borrowed. JMdict also writes it with characters for \textquotedblleft{}ball\textquotedblright{} and \textquotedblleft{}child\textquotedblright{}, which matches the usual account: \emph{tama}, \textquotedblleft{}ball\textquotedblright{}, with \emph{ko}, \textquotedblleft{}child\textquotedblright{}, voiced to \emph{go}.
\end{cousinweb}

\begin{grammarlens}[title={What you've built}]
One more word for this chapter.
\end{grammarlens}

\subsection*{Your turn}
\begin{itemize}
\raggedright
  \item \emph{Say it:} \emph{tamago}
  \item \emph{Say it:} \emph{tamago}, once more
  \item \emph{Say it:} say \emph{tamago}
  \item \emph{Recall:} say the Japanese for a fish, then the Japanese for meat, then say \emph{tamago} again
\end{itemize}

\begin{tcolorbox}[breakable,colback=teal!4,colframe=teal!35!black,title={Before you move on}]
What does \ja{たまご} mean? (\textquotedblleft{}An egg\textquotedblright{}.) Say it once more.
\end{tcolorbox}
\section[\texorpdfstring{\ja{やさい} (yasai)}{やさい (yasai)}]{\ja{やさい} (yasai) — vegetables}

\label{lesson:JA-C32-yasai}



\begin{quote}
Before the new one: say the Japanese for meat, then the Japanese for an egg.
\end{quote}

\subsection*{You'll want to know: \ja{やさい}}
\textbf{\ja{やさい}} — \emph{yasai} — \textquotedblleft{}vegetables\textquotedblright{}.

\begin{cousinweb}
\textbf{\ja{やさい}} (\emph{yasai}) is Sino-Japanese (kango): two Chinese characters read with their on'yomi, readings borrowed from Chinese. They are \emph{ya}, \textquotedblleft{}field\textquotedblright{} (Middle Chinese \emph{yaeX}) and \emph{sai}, \textquotedblleft{}vegetables\textquotedblright{} (\emph{tshojH}) (Unihan, KANJIDIC2; Middle Chinese after Baxter–Sagart).
\end{cousinweb}

\begin{grammarlens}[title={What you've built}]
One more word for this chapter.
\end{grammarlens}

\subsection*{Your turn}
\begin{itemize}
\raggedright
  \item \emph{Say it:} \emph{yasai}
  \item \emph{Say it:} \emph{yasai}, once more
  \item \emph{Say it:} say \emph{yasai}
  \item \emph{Recall:} say the Japanese for meat, then the Japanese for an egg, then say \emph{yasai} again
\end{itemize}

\begin{tcolorbox}[breakable,colback=teal!4,colframe=teal!35!black,title={Before you move on}]
What does \ja{やさい} mean? (\textquotedblleft{}Vegetables\textquotedblright{}.) Say it once more.
\end{tcolorbox}
\section[\texorpdfstring{\ja{くだもの} (kudamono)}{くだもの (kudamono)}]{\ja{くだもの} (kudamono) — fruit}

\label{lesson:JA-C32-kudamono}



\begin{quote}
Before the new one: say the Japanese for an egg, then the Japanese for vegetables.
\end{quote}

\subsection*{You'll want to know: \ja{くだもの}}
\textbf{\ja{くだもの}} — \emph{kudamono} — \textquotedblleft{}fruit\textquotedblright{}.

\begin{cousinweb}
\textbf{\ja{くだもの}} (\emph{kudamono}) is a native word, inherited rather than borrowed. The usual account takes it apart as \emph{ku}, an old form of \emph{ki}, \textquotedblleft{}tree\textquotedblright{}, an old linking \emph{da}, and \emph{mono}, \textquotedblleft{}thing\textquotedblright{}: \textquotedblleft{}tree things\textquotedblright{}. JMdict's kanji, \textquotedblleft{}fruit\textquotedblright{} and \textquotedblleft{}thing\textquotedblright{}, are read for their meaning.
\end{cousinweb}

\begin{grammarlens}[title={What you've built}]
One more word for this chapter.
\end{grammarlens}

\subsection*{Your turn}
\begin{itemize}
\raggedright
  \item \emph{Say it:} \emph{kudamono}
  \item \emph{Say it:} \emph{kudamono}, once more
  \item \emph{Say it:} say \emph{kudamono}
  \item \emph{Recall:} say the Japanese for an egg, then the Japanese for vegetables, then say \emph{kudamono} again
\end{itemize}

\begin{tcolorbox}[breakable,colback=teal!4,colframe=teal!35!black,title={Before you move on}]
What does \ja{くだもの} mean? (\textquotedblleft{}Fruit\textquotedblright{}.) Say it once more.
\end{tcolorbox}
